% GENERATED by scripts/build_data.py; edit JSON inputs, not this file.
% Requires amsmath, booktabs, array, tabularx and cite in the parent preamble.
\ifdefined\TOneTableBox\else\newsavebox{\TOneTableBox}\fi
\begin{table*}[!t]
\centering
\caption{Literature-grounded streaming and resident-write capabilities of representative CIM implementations}
\label{tab:t1_capabilities}
\begingroup
\setlength{\tabcolsep}{3pt}
\renewcommand{\arraystretch}{1.10}
\sbox{\TOneTableBox}{\begin{minipage}{\textwidth}\fontsize{8.5}{10.2}\selectfont
\begin{tabularx}{\textwidth}{@{}>{\raggedright\arraybackslash}p{31mm}>{\raggedright\arraybackslash}p{40mm}rrr>{\raggedright\arraybackslash}X@{}}
\toprule
Technology / implementation & Mode; logical $n_{\rm out}\!\times n_{\rm in}$; input/weight & $\rho$ & $\tau$ & $\mathrm{RI}^{*}$ & Evidence / key condition \\
 & & \multicolumn{2}{c}{[GB/s]} & & \\
\midrule
SRAM / Jia et al.~\cite{t1_jia2022} & ACIM; $64\times1152$; 4b/4b & $2.9$ & \textemdash & \textemdash & M/C; 4 input phases; 20 MS/s ADCs. \\
SRAM / D6CIM~\cite{t1_d6cim2023} & DCIM; $16\times128$; 8b/8b & $0.72$ & \textemdash & \textemdash & M/C; 64 clocks/MVM at 360 MHz. \\
DRAM / Ambit~\cite{t1_ambit2017} & In-array Boolean row logic & \textemdash & \textemdash & \textemdash & S/D; 49 ns AAP; no full MVM schedule. \\
Gain-cell eDRAM / Xie~\cite{t1_gaincell2022} & ACIM; $1\times27$; 8b/8b & $1.4$ & \textemdash & \textemdash & M/C; one MAV group, 4 phases / 20 ns. \\
3D NAND / Shim--Yu~\cite{t1_nand2021} & ACIM model; $16\times4608$; 8b/8b & \textemdash & \textemdash & \textemdash & S/D; 303 ns WL setup only; verify untimed. \\
2D NOR Flash / Guo~\cite{t1_nor2017,t1_nor_tuning2016} & Analog cascade; $784\to64\to10$ & \textemdash & \textemdash & \textemdash & M/D; $<1\,\mu$s classification; preload excluded. \\
RRAM / NeuRRAM~\cite{t1_neurram2022} & ACIM; $256\times128$; 4b/analog & $0.016$ & \textemdash & \textemdash & M/C; 3.9 $\mu$s; 56 $\mu$s/cell write is projected. \\
MRAM / Jung~\cite{t1_jung2022} & ACIM; $64\times64$; bipolar 1b/1b & $0.089$ & $0.044$ & $2$ & M/C; row write: 64 weights in 2 clocks. \\
PCM / 64-core chip~\cite{t1_pcm64} & ACIM; $256\times256$; 8b/analog & $0.49$ & \textemdash & \textemdash & S/C; measured-chip RTL, 4 phases / 520 ns. \\
FeRAM / C2FeRAM~\cite{t1_c2feram2023} & Proposed ACIM; discrete 2T2C & \textemdash & \textemdash & \textemdash & D/D; device tests; no complete timed macro. \\
3D FeNOR / Zhou 2025~\cite{t1_zhou2025} & Device multiply; binary state & \textemdash & \textemdash & \textemdash & D/D; +4.25/-3.25 V, 50 ns; RAW delay $<100$ ns. \\
Vertical FeFET / Zhou 2026~\cite{t1_zhou2026} & Memory array; binary state & \textemdash & \textemdash & \textemdash & D/D; $\pm2$ V, 20 ns; 1 Kib patterns verified. \\
\bottomrule
\end{tabularx}
\par\vspace{5pt}
\textit{Scope:} Native local service boundaries; one core/group unless specified. Packed logical input bytes are counted once per complete stated-precision evaluation; physical bit slices and complementary cells are not extra payload. Analog outputs retain the reported quantization/error contract. $\mathrm{RI}^{*}=\rho/\tau$ only for a compatible pair; GB/s is decimal. These configurations are not normalized for process, area, capacity or power.
\par\vspace{2pt}\textit{Evidence:} First code: M, measured chip; D, measured device; S, simulation. Second code: D, directly reported parameter; C, converted capability. PCM timing is RTL-derived on a measured-chip architecture. \textemdash\ means insufficient evidence for this service boundary, not impossibility. AAP is a DRAM row-operation primitive; RAW means read-after-write. FeNOR switching pulses and RAW delays do not specify CIM throughput. Full conditions, write/verify evidence and unavailable timing are retained in the supporting records.
\end{minipage}}
\typeout{T1_BODY_WIDTH_PT=\the\wd\TOneTableBox}
\typeout{T1_BODY_HEIGHT_PT=\the\ht\TOneTableBox}
\typeout{T1_BODY_DEPTH_PT=\the\dp\TOneTableBox}
\usebox{\TOneTableBox}
\endgroup
\end{table*}
